\appendix

\section{Curve construction and fit diagnostics}
\label{app:curve}

The U.S. Treasury publishes par yields at eleven tenors. Zero rates are recovered by
forward substitution: with nodes $z(\tau_1), \dots, z(\tau_{k-1})$ already determined,
$z(\tau_k)$ is the rate that makes a par bond maturing at $\tau_k$, paying the
published coupon semi-annually, price to one hundred. Coupon dates falling between
existing nodes are discounted at a rate interpolated linearly between the last fixed
node and the trial value of $z(\tau_k)$, so the construction is self-consistent
rather than dependent on an interpolation applied after the fact. The residual on
re-pricing each par bond from the completed curve is below one basis point of yield at
every tenor and both dates.

The bootstrapped zeros are then fitted with the Svensson form, in which the zero rate
at tenor $\tau$ is

\begin{equation}
z(\tau) = \beta_0
 + \beta_1 \frac{1 - e^{-\tau/\lambda_1}}{\tau/\lambda_1}
 + \beta_2 \left( \frac{1 - e^{-\tau/\lambda_1}}{\tau/\lambda_1} - e^{-\tau/\lambda_1} \right)
 + \beta_3 \left( \frac{1 - e^{-\tau/\lambda_2}}{\tau/\lambda_2} - e^{-\tau/\lambda_2} \right).
\label{eq:nss}
\end{equation}

Fitting minimises error in \emph{par} yield rather than in the bootstrapped zeros,
because the par yields are what Treasury actually published and the zeros are already
a construction. The objective is non-convex and admits degenerate optima --- an
unconstrained fit at the end of 2021 returns a level parameter of $-36\%$ with the
second curvature term pinned at its bound, fitting the observed tenors acceptably
while implying negative forward rates beyond thirty years. Fitting is therefore
multi-start, with $\beta_0$ bounded to $[0, 0.15]$, and candidate solutions are
rejected if the implied forward curve goes negative anywhere in the fitted range.

Table~\ref{tab:residuals} reports the fit at both valuation dates.

\begin{table}[htbp]
\centering
\caption{Par-yield fit at both valuation dates, in per cent, with residuals in basis
points. Residual is fitted less observed. Root-mean-square error is 3.88 basis
points at the end of 2021 and 5.39 at the end of 2022.}
\label{tab:residuals}
\small
\begin{tabular}{lrrrrrr}
\toprule
 & \multicolumn{3}{c}{31 December 2021} & \multicolumn{3}{c}{31 December 2022} \\
\cmidrule(lr){2-4} \cmidrule(lr){5-7}
Tenor & Observed & Fitted & Residual & Observed & Fitted & Residual \\
\midrule
1m & 0.06 & 0.05 & $-1.37$ & 4.12 & 4.13 & $+0.52$ \\
3m & 0.06 & 0.09 & $+2.86$ & 4.42 & 4.44 & $+1.96$ \\
6m & 0.19 & 0.18 & $-1.34$ & 4.76 & 4.69 & $-6.63$ \\
1y & 0.39 & 0.38 & $-0.76$ & 4.73 & 4.79 & $+6.14$ \\
2y & 0.73 & 0.74 & $+0.82$ & 4.41 & 4.42 & $+0.77$ \\
3y & 0.97 & 0.98 & $+0.55$ & 4.22 & 4.17 & $-4.84$ \\
5y & 1.26 & 1.25 & $-1.12$ & 3.99 & 3.98 & $-0.90$ \\
7y & 1.44 & 1.41 & $-2.55$ & 3.96 & 3.95 & $-1.12$ \\
10y & 1.52 & 1.58 & $+6.38$ & 3.88 & 3.97 & $+9.02$ \\
20y & 1.94 & 1.85 & $-8.70$ & 4.14 & 4.04 & $-10.05$ \\
30y & 1.90 & 1.95 & $+5.25$ & 3.97 & 4.02 & $+5.00$ \\
\bottomrule
\end{tabular}
\end{table}

The residuals are largest at the twenty-year point at both dates, which is where the
published curve has a kink the four-factor form cannot follow: the thirty-year par
yield is below the twenty-year at both dates, and a Svensson curve smooths through
that rather than reproducing it. Since both banks held material positions in the
over-ten-years bucket, this is the part of the fit that matters most for the result,
and the error it introduces is of the same order as the 8.7 to 10 basis points
reported rather than negligible.

\section{Short-rate calibration and pricing validation}
\label{app:pricing}

Under the Hull--White dynamics of equation~\eqref{eq:hw}, a zero-coupon bond has the
closed form

\begin{equation}
P(t,T) = A(t,T)\, e^{-B(t,T)\, r(t)},
\qquad
B(t,T) = \frac{1 - e^{-a(T-t)}}{a},
\label{eq:zcb}
\end{equation}

with $A(t,T)$ determined by the initial curve and the volatility. The drift function
$\theta(t)$ is fitted analytically to the forward curve of
Appendix~\ref{app:curve}, so the model reproduces the initial term structure by
construction and no fitting error is introduced at $t=0$.

Simulation works on the decomposition $r(t) = x(t) + \alpha(t)$ described in
Section~\ref{sec:reconstruction}. Because $x$ is an Ornstein--Uhlenbeck process with
known transition density, it can be stepped exactly at any time increment, so the
only discretisation error in a path is in the integral of $r$ accumulated for
discounting, not in the short rate itself.

The validation question is whether simulated discount factors average back to the
curve the model was built on. If they do not, every price computed downstream is
wrong in a way calibration will not reveal, since the option-adjusted spreads would
simply absorb the error. Table~\ref{tab:mc} reports the comparison at the production
path count and seed.

\begin{table}[htbp]
\centering
\caption{Simulated against closed-form discount factors, 20,000 antithetic paths,
monthly steps, $a = 0.0889$ and $\sigma = 61.9$ basis points. Error is in basis
points of continuously compounded yield; the final column expresses the discrepancy
in Monte Carlo standard errors.}
\label{tab:mc}
\small
\begin{tabular}{rrrrrr}
\toprule
Maturity & Simulated & Closed form & Error (bp) & Std.\ error & Error in s.e. \\
\midrule
1  & 0.996188 & 0.996191 & $+0.0294$ & 0.000024 & $-0.120$ \\
5  & 0.939254 & 0.939261 & $+0.0161$ & 0.000224 & $-0.034$ \\
10 & 0.852332 & 0.852340 & $+0.0083$ & 0.000501 & $-0.014$ \\
20 & 0.686029 & 0.685913 & $-0.0840$ & 0.000890 & $+0.130$ \\
30 & 0.549768 & 0.549573 & $-0.1184$ & 0.001060 & $+0.184$ \\
\bottomrule
\end{tabular}
\end{table}

The worst error is 0.12 basis points of yield, at thirty years, and the largest
discrepancy in standard-error units is 0.18. The engine sits well inside sampling
noise at every tenor, and the test suite gates the pipeline at five basis points.
These are the only standard errors in this paper, and they describe the numerical
precision of a simulation whose answer is known analytically; they carry no
information about uncertainty in any economic quantity.

The parameters themselves are estimated from the three-month bill, and
Section~\ref{sec:sensitivity} reports the estimator across windows and sampling
frequencies. Two points from that table belong here. The specification used --- monthly
observations from 1990 --- gives a mean-reversion half-life of 7.80 years, and a
configuration guard rejects any fit implying a half-life beyond fifty years, which is
what prevents the degenerate 2003--2022 estimates from entering the model silently.
The honest qualification is that mean reversion is weakly identified under every
specification tried, with a best $R^2$ of 0.0387.

\section{The deposit franchise and the implied half-life}
\label{app:deposits}

Franchise value follows equation~\eqref{eq:deposits}: the book balance less the
present value of the below-market funding benefit, with the balance decaying
exponentially at a stated half-life and the benefit per period equal to the retained
spread $(1-\beta) y_t$ net of a servicing cost. Deposits are split into insured and
uninsured layers with separate half-lives, and the uninsured layer is the one the run
model allows to leave.

The solve in Section~\ref{sec:comparison} fixes $\beta$ at the 0.60 SVB discloses and
the insured half-life at 3.5 times the uninsured one, preserving the ratio of the
conservative case, then searches for the uninsured half-life $h$ at which the modelled
change in economic value at $+100$ basis points equals the disclosed $-13.9\%$. The
search uses Brent's method over $h \in [0.25, 25]$ years, bracketed by a sign change
at the endpoints and converged to a tolerance of $10^{-4}$ years.

A bracketing sign change establishes that a root exists but not that it is unique, so
monotonicity was checked rather than assumed. Evaluating the sensitivity at 38 points
across the search interval gives a function that is strictly increasing throughout,
from $-67.90\%$ at three months to $+16.47\%$ at twenty-four and a half years, with no
change of sign in the first difference anywhere. The root is therefore unique, and
the solver returns $h = 3.5896$ years with the insured half-life following at
12.5636.

The direction has a straightforward reading. Lengthening the assumed half-life
increases the present value of the below-market funding benefit, which raises both
base economic value and the gain the liability side records when rates rise, and that
gain offsets more of the asset-side loss. A bank that assumes its depositors stay
longer reports a smaller exposure, monotonically and without limit over any range a
reviewer would entertain.

Two properties of the solution are worth stating. It is a one-parameter fit to a
one-number target, so it cannot fail to hit that target and the agreement at $+100$
basis points carries no information; this is why the paper treats only the other
shock points as evidence. And it is conditional on everything else in the model being
right, so any bias in the asset valuation is transferred into the implied half-life
rather than showing up as a residual. The asset-side conservatism documented in
Section~\ref{sec:comparison} therefore pushes the implied half-life upward, in the
same direction as the hedging conflation discussed in
Section~\ref{sec:interpretation}. Both work against the paper's argument, and the
argument rests on the width of the sensitivity in Section~\ref{sec:sensitivity}
rather than on the point estimate.
